\section{Conclusions}
\label{sec:conclusion}

This study presented a simulation-grounded multi-agent LLM framework for
traffic optimization in SUMO.  The framework separates proposal generation
from decision making: specialist agents generate signal-control,
corridor-speed, and routing interventions from compact diagnostic evidence,
while deterministic capability validation prevents unsupported actions from
reaching execution.  Every valid single intervention and compatible pairwise
bundle is evaluated through a complete simulation.  Final selection uses a
baseline-frozen bottleneck objective with explicit constraints on global delay,
neighboring-road spillover, throughput, and teleport events, together with an
unmodified no-operation candidate.  An offline Pareto audit further checks the
selected plan against delay, emissions, and mobility outcomes without changing
the optimizer's primary objective.

Across five road networks and three demand seeds, the guarded optimizer reduced
Top-10 bottleneck time loss by $9.18\%\pm10.18\%$ and network-wide delay by
2.78\%.  It selected a feasible improving plan in 86.7\% of cases, compared
with 73.3\% for grounded random search and 80.0\% for the deterministic
heuristic under the same simulation cap.  Feedback-refined candidates appeared
in seven final plans and bundles in four, showing that the second-stage search
often changed the executable decision.  Eleven of the 15 selected plans were
nondominated within their evaluated candidate pools.

The current findings are limited by synthetic demand, three demand seeds, a
delay-oriented proposal prompt, and the absence of a budget-matched conventional
optimizer.  The small Pareto sets further indicate that the present candidate
pool contains limited environmental and mobility diversity.  Future work will
generate candidates using objective-specific prompts and multiple LLM
backends, compare the framework with non-LLM search and established traffic
control baselines, and evaluate calibrated or observed demand.  Extensions to
larger intervention bundles, multiple coordinated controllers, and time-varying
online policies are also needed.  Overall, the results suggest that LLMs are
most useful as adaptable generators of structured traffic interventions, while
executable tools, microscopic simulation, and explicit safety constraints
should retain authority over the final decision.
